\chapter{提示词工程}
\section{什么是提示词工程}
在前面的章节中，我们学习了从感知机到 Transformer 的深度学习技术。当我们训练好一个大型语言模型（如 GPT、Claude 等）后，如何使用它就成为了一个关键问题。
提示词工程（Prompt Engineering）就是研究如何设计和优化输入提示（Prompt），以引导模型生成更准确、更有用的输出。

如果把大语言模型比作一个知识渊博的专家，那么提示词就是你向这位专家提问的方式。
同一个问题，不同的问法可能得到完全不同质量的答案。提示词工程就是学习如何“问对问题”的艺术。

大型语言模型本质上是一个条件概率模型，给定输入序列 $x_1, x_2, \ldots, x_n$，模型预测下一个词 $x_{n+1}$ 的概率分布：

\begin{equation*}
	P(x_{n+1} | x_1, x_2, \ldots, x_n)
\end{equation*}

提示词就是这个条件序列，一个好的提示能够更好地激活模型中相关的知识，引导模型生成更符合期望的输出。
这就像在图书馆找书，如果你能给图书管理员更精确的描述，他就更容易找到你想要的书。
提示词工程不需要修改模型的参数，只需要调整输入，因此它是优化模型表现的一种非常灵活和低成本的方式。

\section{常见的提示方法}

\subsection{零样本提示}
零样本提示（Zero-shot Prompting）是最简单的提示方式，就是直接告诉模型你要什么，不提供任何示例，
模型会根据它预训练时学到的知识直接给出答案。

\begin{verbatim}
提示：将以下英文翻译成中文：
"Hello, how are you?"

输出：你好，最近怎么样？
\end{verbatim}

\subsection{少样本提示}
少样本提示（Few-shot Prompting）是在提示中提供几个示例，让模型学习你想要的格式和模式。这就像教新员工时，你会先做几个示范，然后让他们模仿。

\begin{verbatim}
提示：
英文：I love this movie -> 情感：正面
英文：This is terrible -> 情感：负面
英文：It's okay -> 情感：中性
英文：What an amazing experience -> 情感：

输出：正面
\end{verbatim}

在这个例子中，我们提供了三个示例，模型就能理解我们想要进行情感分类任务，并按照相同的格式输出。
少样本提示之所以有效，是因为大型语言模型具有强大的上下文学习能力（In-context Learning）。模型虽然没有在训练时专门学习这个任务，
但它能够根据提示中的示例，推断出任务的模式和期望的输出格式。

\subsection{思维链提示}

思维链提示（Chain-of-Thought Prompting，CoT）是一种让模型“逐步思考”的技术。
对于复杂的推理任务，直接要求答案往往效果不好，但如果让模型展示思考过程，准确率会大幅提升。

这就像做数学题时，老师会要求你写出解题步骤，而不仅仅是答案。写出步骤不仅能帮助你理清思路，也更容易发现错误。

\begin{verbatim}
提示：
问题：小明有 5 个苹果，他给了小红 2 个，又买了 3 个，现在他有几个苹果？
让我们一步一步思考。

输出：
1. 小明最初有 5 个苹果
2. 他给了小红 2 个，剩下 5 - 2 = 3 个
3. 他又买了 3 个，所以 3 + 3 = 6 个
4. 答案是：小明现在有 6 个苹果
\end{verbatim}

关键在于“让我们一步一步思考”这句话，它引导模型展示推理过程，而不是直接猜测答案。

\section{提示词工程的实践技巧}

在实践中，有几个常用的技巧可以显著提升提示的质量。

\textbf{明确具体。}模糊的提示往往得到模糊的回答。好的提示应该明确说明你的需求：

\begin{verbatim}
差的提示：写点关于Python的内容

好的提示：用Python写一个函数，输入是一个整数列表，
输出是列表中的最大值。要求使用for循环实现，
并添加类型注解和文档字符串。
\end{verbatim}

\textbf{指定格式。}如果你需要特定格式的输出，明确告诉模型：

\begin{verbatim}
提示：列出三个Python的优点，用JSON格式输出，
每个优点包含"title"和"description"字段。
\end{verbatim}

\textbf{角色扮演。}让模型扮演特定角色，可以得到更专业的回答：

\begin{verbatim}
提示：你是一位有10年经验的高级Python开发工程师。
请review以下代码，指出潜在的问题和改进建议：
[代码内容]
\end{verbatim}

\textbf{分解复杂任务。}对于复杂的任务，将其分解为多个简单的步骤：

\begin{verbatim}
提示：
任务：分析一篇技术文章并生成摘要

步骤1：提取文章的三个核心观点
步骤2：为每个观点写一句话总结
步骤3：将三个总结组合成一段100字以内的摘要

文章内容：[文章内容]
\end{verbatim}
